\documentclass{article}
\usepackage[UTF8]{ctex}
\usepackage{amsmath,amssymb,amsfonts}
\usepackage{enumitem}
\usepackage{geometry}
\geometry{a4paper, margin=1in}
% === 错题本排版优化（不参与题目标识） ===
\setlength{\parindent}{0pt}
\setlength{\parskip}{0.35em plus 0.10em minus 0.08em}
\setlength{\emergencystretch}{3em}
\setlist{itemsep=0.25em, topsep=0.40em, parsep=0pt}
\allowdisplaybreaks[2]
% === 错题本排版优化结束 ===
\begin{document}
\section*{西南交通大学2026年高等代数真题}
\section*{163．西南交通大学2026年高等代数试卷（回忆版）}
（注．题目来自考生回忆，部分表述或数据可能有误）

一．填空题．\\

\subsection{题目1}
1．设 \(\sigma(\alpha)=\alpha-k(\alpha, \varepsilon) \varepsilon, k\) 为常数，\(\varepsilon\) 为单位列向量，若 \(\sigma\) 是正交变换，则 \(k=\) \(\_\_\_\_\)．

\subsection{题目2}
2．设 \(\mathbb{Q}(\sqrt{2}, \sqrt{3})\) 是数域 \(\mathbb{Q}\) 上的最小线性空间，求其维数 \(\_\_\_\_\)．\\
\(\_\_\_\_\)

\subsection{题目3}
3．设 \(A\) 为3阶正定阵，在条件 \(x_{1}^{2}+x_{2}^{2}+x_{3}^{2}=1\) 下，二次型 \(\left(x_{1}, x_{2}, x_{3}\right) A\left(x_{1}, x_{2}, x_{3}\right)^{\mathrm{T}}\) 的最大值，最小值分别为 4,1 ，且 \(|A|=8\) ，则 \(\operatorname{tr}(A)=\) \(\_\_\_\_\)．

\subsection{题目4}
4．设多项式 \(f(x)=x^{3}+2 x^{2}-2, g(x)=x^{2}+x-1\) ，若 \(a, b, c\) 为 \(f(x)\) 的根，\(g(a), g(b), g(c)\) 是 \(h(x)\)的根，则 \(h(x)=\) \(\_\_\_\_\)．

\subsection{题目5}
5．以矩阵 \(\left(\begin{array}{cccc}1 & 0 & a & b \\ 0 & 1 & c & 0 \\ a & c & 0 & -e \\ b & a & e & 0\end{array}\right)\) 为系数阵的方程组若以 \(x_{3}, x_{4}\) 为自由未知量，求 \(a, b, c, d, e\) 满足的关系\\
\(\_\_\_\_\)．

二．解答题。\\

\subsection{题目6}
1．设 \(A=\alpha \beta^{\mathrm{T}}\) ，且 \(\alpha^{\mathrm{T}} \beta=0\) ，求 \(A\) 的特征值和特征子空间的一组基。

\subsection{题目7}
2．设 3 阶实对称矩阵 \(A\) 有特征值 \(1,2,2\) ，特征值1对应特征向量为 \((1,1,1)^{\mathrm{T}}\) ，求 \(A^{3}\) ：

\subsection{题目8}
3．设 \(A=\left(a_{i j}\right), B=\left(b_{i j}\right) \in M_{2}(\mathbb{R})\) ，定义内积 \((A, B)=\sum_{i=1}^{2} \sum_{j=1}^{2} a_{i j} b_{i j}\) 。取

\[
A_{1}=\left(\begin{array}{ll}
1 & 0 \\
0 & 1
\end{array}\right), A_{2}=\left(\begin{array}{ll}
0 & 1 \\
1 & 1
\end{array}\right) .
\]

令 \(W=L\left(A_{1}, A_{2}\right)\) ，求 \(W^{\perp}\) 的一组标准正交基。

\subsection{题目9}
4．设矩阵

\[
A=\left(\begin{array}{ccc}
1 & -2 & 6 \\
-1 & 0 & 3 \\
-1 & -1 & 4
\end{array}\right)
\]

\begin{itemize}
  \item[（1）] 求 \(A\) 的 Jordan 标准形．
  \item[（2）] 求 \(B, C\) 使得 \(A=B+C\) ，其中 \(B\) 是幂零阵，\(C\) 为可对角化的矩阵．
\end{itemize}

\subsection{题目10}
5．设 \(f(x)=x^{5}+m x+1\) ，求整数 \(m\) 的值使得 \(f(x)\) 在 \(\mathbb{Q}\) 上可约．

\subsection{题目11}
6．设 \(V\) 为 \(n\) 维线性空间，\(V_{1}\) 为 \(V\) 的子空间，若 \(\operatorname{dim} V_{1} \geq \frac{n}{2}\) ，证明：存在 \(V\) 的子空间 \(W_{1}, W_{2}\) ，使得\(V=V_{1} \oplus W_{1}=V_{1} \oplus W_{2}\) ，且 \(W_{1} \cap W_{2}=\{0\}\) ．

\subsection{题目12}
7．设 \(A, B\) 为同阶正交阵。

\begin{itemize}
  \item[（1）] 若 \(|A|+|B|=0\) ，证明：\(|A+B|=0\) ．
  \item[（2）] 若 \(n\) 为奇数，证明：\(|(A-B)(A+B)|=0\) ．
\end{itemize}

\subsection{题目13}
8．设 \(n\) 阶矩阵 \(A\) 有 \(n\) 个不同的特征值，\(n\) 阶矩阵 \(B\) 满足 \(A B=B A\) ，证明：存在多项式 \(f(x)\) ，使得\(B=f(A)\) ．

\subsection{题目14}
9．设 \(A\) 是 \(n\) 阶半正定矩阵，\(B\) 是 \(n\) 阶正定矩阵，证明：\(|A+B| \geq|B|\) ，且等号成立当且仅当 \(A=O\) 。
\end{document}
